% ============================================================
%  sections/experiments.tex  —  Experiments
% ============================================================
\section{Experiments}
\label{sec:experiments}

We evaluate SceneFactory along three axes:
simulation throughput (\S\ref{sec:throughput}),
end-to-end navigation performance on held-out Waymo scenes (\S\ref{sec:navigation_perf}),
and qualitative trajectory analysis (\S\ref{sec:qualitative}).
All experiments use a single shared navigation policy trained with PPO; the contribution of this work is the simulation platform itself, so we report policy performance as an end-to-end validation of the pipeline rather than as a claim of state-of-the-art driving.

% ---------------------------------------------------------------
\subsection{Experimental Setup}
\label{sec:exp_setup}
% ---------------------------------------------------------------

\paragraph{Hardware and software.}
All training and evaluation runs execute on a single NVIDIA RTX PRO 6000 Blackwell GPU (96\,GB VRAM).
The software stack consists of
NVIDIA Isaac Sim~\chattodo{version number},
Isaac Lab~\chattodo{version number}~\cite{nvidia_isaac_2025},
RSL-RL~\cite{schwarke_rsl-rl_2025} for PPO,
PyTorch~\chattodo{version}, and
CUDA~\chattodo{version}.
Physics runs at 120\,Hz (time step $\Delta t = 1/120$\,s);
the control policy acts at 30\,Hz (decimation factor 4).

\paragraph{Scene curation.}
From approximately 1{,}000 Waymo Open Motion Dataset scenarios~\cite{Kan_2024_icra},
we extract candidate road topologies through the pipeline described in \S\ref{sec:data_processing}.
A human operator first visually inspects each candidate scene with a curation wizard that renders lane-center polylines, road-edge polylines, and origin-destination pairs.
The operator accepts scenes that have plausible road geometry and rejects those with degenerate lane structure, multi-level overlaps from z-flattening, or unreachable goals (see \S\ref{sec:sharp_edges} for rejection criteria).
This produces \textbf{64 training scenes} and a disjoint set of \textbf{64 held-out test scenes}.
Within each accepted scene, an automated filter then retains only those agents whose start and goal positions both lie within the 400\,m $\times$ 400\,m world bounding box, are separated by at least 3.0\,m, and have valid coordinates.
Up to 16 qualifying agents are extracted per scene; scenes with fewer valid agents leave the remaining agent slots inactive.

\paragraph{Training configuration.}
Each training scene is replicated four times to fill $64 \times 4 = 256$ parallel environments, each hosting up to 16 agent slots, for a nominal capacity of 4{,}096 agents.
In practice, many Waymo scenes contain fewer than 16 qualifying agents after the spawn filter described below, so slots without a valid spawn remain inactive throughout the episode.
Worlds are assigned to environments by random shuffling (seed~42), so every PPO update samples transitions from all 64 unique road topologies.
Episodes last 50\,s (1{,}500 control steps).
Upon reaching the goal (within 3.0\,m) or triggering a termination condition (collision, off-road crash, timeout), an agent is parked in place with zero velocity; the remaining agents continue until every agent in the environment has terminated or the episode times out, at which point the entire world resets with fresh start-goal assignments for all agents.

\paragraph{PPO hyperparameters.}
Training uses PPO with an adaptive learning rate (initial $2.09{\times}10^{-4}$, KL target 0.01), rollout horizon of 128 steps, 5 gradient epochs over 16 mini-batches, $\gamma{=}0.99$, $\lambda{=}0.98$, and clip $\epsilon{=}0.15$.
The total batch size per update is $4{,}096 \times 128 = 524{,}288$ transitions.
The full hyperparameter table is in Appendix~\ref{app:ppo_hparams}.

\paragraph{Evaluation protocol.}
We evaluate on the 64 held-out test scenes, each loaded into a separate environment with up to 16 agent slots.
Across the 64 test worlds, 624 of the 1{,}024 nominal slots contain a valid spawn (the remainder correspond to scenes with fewer than 16 qualifying agents).
The policy runs in deterministic mode (mean of the action distribution).
Evaluation uses \emph{invincible mode}: collisions and off-road events are logged but do not terminate the agent, allowing it to continue toward the goal so that we can measure the unconditional goal-reach rate.
An agent is counted as successful if it reaches within 3.0\,m of its assigned goal before the 50\,s episode deadline.

% ---------------------------------------------------------------
\subsection{Throughput}
\label{sec:throughput}
% ---------------------------------------------------------------

A core engineering contribution of SceneFactory is throughput: training with full rigid-body dynamics at a rate that makes large-scale PPO practical on a single GPU.
To quantify this, we compare SceneFactory against its direct predecessor, a pipeline built on the same Isaac Sim simulator and PhysX 5 solver but using the built-in Vehicle Wizard vehicle model and Stable-Baselines3 for PPO (hereafter ``Baseline'').
Both pipelines load the same Waymo scene, use 16 agents per world, identical observation features (350 road points, 24 vehicle neighbors with TTC), and run on the same RTX PRO 6000 Blackwell GPU.
The only differences are the vehicle representation, the environment wrapper architecture, and the RL framework.

\paragraph{Metric.}
We report \emph{controlled agent steps per second} (CASPS): the number of policy-controlled agents that complete one 30\,Hz control tick per wall-clock second.
Both pipelines log CASPS every PPO iteration via embedded sub-phase timers;
we discard the first 50 logged iterations as warmup and report the mean and standard deviation over the remaining steady-state samples.

\paragraph{Scaling experiment.}
We select a single Waymo scene with 28 qualifying agents (ensuring all 16 slots are filled in every world) and replicate it at four scales: 32, 64, 128, and 256 parallel worlds, yielding 512 to 4{,}096 agent slots.
Each configuration trains for a short benchmark window: 500\,K total timesteps for the Baseline, 20 PPO iterations for SceneFactory.
Figure~\ref{fig:throughput_scaling} plots CASPS as a function of total agent slots for both pipelines.
The Baseline throughput is essentially flat at $\sim$152--174 CASPS across all world counts, confirming that per-agent Python loops---not the PhysX solver---are the bottleneck.
SceneFactory scales near-linearly, reaching 19{,}250 CASPS at 256 worlds (4{,}096 agent slots): a \textbf{127$\times$ speedup} over the Baseline at the same scale.
Even at the smallest configuration (32 worlds), SceneFactory is 22$\times$ faster.
Neither pipeline exceeded GPU memory at 256 worlds on the RTX PRO 6000 Blackwell (96\,GB VRAM).

\begin{table}[t]
\centering
\caption{Throughput scaling (CASPS, mean $\pm$ std over steady-state iterations) on a single RTX PRO 6000 Blackwell GPU.
Both pipelines use the same Waymo scene, the same PhysX 5 rigid-body solver, and 16 agents per world.}
\label{tab:throughput_scaling}
\small
\begin{tabular}{rrcc}
\toprule
Worlds & Agent slots & Baseline & SceneFactory \\
\midrule
 32 &   512 & $174 \pm 27$ & $3{,}870 \pm 66$ \\
 64 & 1{,}024 & $159 \pm 29$ & $7{,}173 \pm 114$ \\
128 & 2{,}048 & $164 \pm 1$ & $12{,}225 \pm 1{,}520$ \\
256 & 4{,}096 & $152 \pm 0$ & $19{,}250 \pm 2{,}984$ \\
\bottomrule
\end{tabular}
\end{table}

\paragraph{Why the speedup.}
The speedup arises from three architectural changes: (1)~GPU-resident joint-state tensors that replace per-agent Python loops with batched \texttt{torch} operations, (2)~cloned PhysX solver islands that enable parallel broad-phase/narrow-phase dispatch, and (3)~an on-device training loop (RSL-RL) that eliminates CPU--GPU data transfers.
A detailed per-phase timing breakdown (physics, observations, rewards, actions) is provided in Appendix~\ref{app:throughput}.

\paragraph{Comparison with kinematic simulators.}
Table~\ref{tab:cross_simulator} places SceneFactory in the context of published throughput numbers from kinematic driving simulators.
GPUDrive~\cite{kazemkhani_gpudrive_2025} (ICLR 2025) reports a peak of 2.3\,M agent-steps per second (ASPS) on an A100 GPU using a kinematic bicycle model.
When restricted to \emph{controlled} agents---i.e., those with non-trivial goals---the effective rate drops to $\sim$200\,K CASPS due to the high fraction of parked or near-goal vehicles in randomly sampled Waymo scenes (mean $\sim$10.8 controllable agents per scene out of up to 128 agent slots).
During end-to-end PPO training the authors report 200\,K--500\,K CASPS depending on hardware and observation configuration.
Nocturne~\cite{vinitsky_nocturne_2022} (NeurIPS 2022), a CPU-based C++ simulator with a kinematic model, achieves $\sim$15\,K ASPS on 16 CPU cores.
Waymax~\cite{gulino_waymax_2023} (NeurIPS 2024), a JAX-based simulator, supports GPU acceleration but encountered out-of-memory errors beyond 16 parallel environments in the GPUDrive authors' benchmarks, so we do not list a peak throughput.

SceneFactory operates at a lower absolute throughput than GPUDrive because each agent is a 10-DOF articulated rigid body simulated by the PhysX 5 solver at 120\,Hz (four physics sub-steps per 30\,Hz control tick).
The comparison is not apples-to-apples: the physics fidelity that makes SceneFactory useful for sim-to-real transfer---tire slip, suspension dynamics, inertial coupling---is precisely what makes it slower than a kinematic model.
Nevertheless, at 19.3\,K CASPS SceneFactory already exceeds Nocturne's throughput while providing full rigid-body dynamics, and it operates within an order of magnitude of GPUDrive's training-time CASPS despite the fundamentally heavier physics workload.
The relevant same-hardware, same-physics comparison is against the Baseline pipeline on the same RTX PRO 6000, where SceneFactory achieves a \textbf{127$\times$} speedup (Table~\ref{tab:throughput_scaling}).

\begin{table}[t]
\centering
\caption{Cross-simulator throughput comparison using published numbers.
All kinematic simulators use simplified vehicle dynamics (no rigid-body contact, no tire model); SceneFactory is the only entry with full PhysX 5 rigid-body physics.
ASPS = agent steps per second (all agents, including parked); CASPS = controlled agent steps per second (policy-active agents only).
GPUDrive and Nocturne numbers are taken from~\cite{kazemkhani_gpudrive_2025}; SceneFactory numbers are measured on our hardware.}
\label{tab:cross_simulator}
\small
\begin{tabular}{llllr}
\toprule
Simulator & Physics model & GPU & Metric & Throughput \\
\midrule
GPUDrive~\cite{kazemkhani_gpudrive_2025}
  & Kinematic bicycle & A100 80\,GB & ASPS & 2{,}300{,}000 \\
GPUDrive (training)
  & Kinematic bicycle & A100 80\,GB & CASPS & 200K--500K \\
Nocturne~\cite{vinitsky_nocturne_2022}
  & Kinematic & CPU 16 cores & ASPS & 15{,}000 \\
\midrule
\textbf{SceneFactory (ours)}
  & \textbf{10-DOF rigid body} & \textbf{RTX PRO 6000} & \textbf{CASPS} & \textbf{19{,}250} \\
\bottomrule
\end{tabular}
\end{table}

% ---------------------------------------------------------------
\subsection{Navigation Performance}
\label{sec:navigation_perf}
% ---------------------------------------------------------------

We train a single shared navigation policy across all 4{,}096 agent slots (256 worlds $\times$ 16 agents) using PPO and select the best checkpoint by validation goal-reach rate.
The selected checkpoint (model\_975) corresponds to approximately $975 \times 128 \times 4{,}096 \approx 511\text{M}$ environment steps.

\paragraph{Evaluation protocol.}
All evaluations use \emph{invincible mode}: collision, lane-boundary, and off-road events are logged but do not terminate the agent, so that every spawned vehicle is observed until it either reaches its goal (within 3.0\,m) or the 50\,s episode budget expires.
This protocol allows us to measure goal-reach rates without survivorship bias.

\paragraph{Generalization study.}
To assess how well the learned policy generalizes across scenes and routing targets, we evaluate the same checkpoint under three conditions that progressively decouple the test setting from the training distribution:
%
\begin{enumerate}[leftmargin=*, itemsep=1pt]
  \item \textbf{Test scenes, original OD.}
        64 held-out Waymo scenes with the original origin--destination pairs extracted from the dataset.
  \item \textbf{Test scenes, random OD.}
        The same 64 held-out scenes, but each agent's destination is re-sampled uniformly along reachable road segments at 20--60\,m travel distance.
  \item \textbf{Train scenes, random OD.}
        The 64 unique training scenes, each with 16 randomly sampled origin--destination pairs (20--60\,m travel distance), for a total of 1{,}024 agents---matching the test-scene agent count.
\end{enumerate}

Table~\ref{tab:generalization} reports the results.

\begin{table}[t]
\centering
\caption{Generalization evaluation of the navigation policy (model\_975).
All runs use invincible mode with a 50\,s episode budget and 3.0\,m goal threshold.
``Original OD'' denotes destination points extracted from the Waymo dataset;
``Random OD'' denotes destinations re-sampled uniformly at 20--60\,m travel distance along reachable road segments.}
\label{tab:generalization}
\small
\begin{tabular}{llccccc}
\toprule
Scenes & OD & Agents & Success $\uparrow$ & Collision $\downarrow$ & Lane-Forbidden $\downarrow$ & $\bar{d}_\text{goal}$ (m) $\downarrow$ \\
\midrule
Test (64)   & Original & 624   & 71.6\% & 6.3\% & 9.0\% & 10.1 \\
Test (64)   & Random   & 1{,}024 & 92.8\% & 4.5\% & 6.9\% & 3.7 \\
Train (64)  & Random   & 1{,}024 & 92.5\% & 3.5\% & 7.6\% & 3.7 \\
\bottomrule
\end{tabular}
\end{table}

On held-out test scenes with original destinations, the policy reaches its goal in 71.6\% of cases (447 of 624 agents) with a mean residual distance of 10.1\,m for non-successful agents.
When the same test scenes are paired with randomly sampled destinations (Condition~2), the success rate jumps to 92.8\% (950 of 1{,}024 agents), nearly matching the 92.5\% achieved on training scenes with the same random-OD protocol (Condition~3).
The near-identical performance between Conditions~2 and~3 is the central finding: \emph{the policy generalizes to entirely novel road geometries as well as it performs on familiar ones}, when the routing difficulty is controlled.
This rules out scene memorization and confirms that the policy has learned a general goal-conditioned navigation strategy.

The 21-percentage-point gap between Conditions~1 and~2 (71.6\% vs.\ 92.8\%), which share the same test scenes but differ only in OD pairs, reveals that the original Waymo destinations are substantially harder than the 20--60\,m random-OD protocol.
The Waymo-extracted routes include long-range multi-turn trajectories, goals near road boundaries, and occasionally ambiguous lane assignments that the random sampler avoids by construction.
This suggests that further gains on the original OD setting are achievable through longer episodes, curriculum over route difficulty, or improved goal-approach behavior, rather than additional scene diversity.

The collision rates remain low across all conditions (3.5--6.3\%), and are inflated by invincible mode, which allows post-collision agents to continue driving and potentially trigger additional collision events; under normal termination semantics each collision would end the episode immediately.
Notably, the off-road crash rate is higher on test scenes (18.7\%) than on training scenes (14.8\%) under the same random-OD protocol, indicating that novel road geometries do cause more lane departures---but the policy recovers and still reaches the goal at nearly the same rate.

\paragraph{Comparison with kinematic simulators.}
GPUDrive reports near-perfect goal-reach rates ($>$98\%) on Waymo scenes using a kinematic bicycle model with no rigid-body dynamics.
We attribute the gap primarily to the harder control problem posed by full PhysX dynamics: the policy must learn to manage tire slip, suspension oscillations, and inertial overshoot that do not exist in kinematic models.
The 10-DOF articulated vehicle with effort-based actuation (\S\ref{sec:action_space}) requires the agent to implicitly discover braking distances, cornering limits, and steering lag, whereas a kinematic model executes velocity commands with no delay or slip.
Closing this gap through longer training, refined reward shaping, and curriculum design is active work; the current results demonstrate that the platform successfully trains navigating agents end-to-end under physically realistic dynamics.

\chattodo{Add a training curve figure (mean reward and/or success rate vs.\ training iteration) for the best run.  Even a single curve is informative.  Export from TensorBoard.}

% ---------------------------------------------------------------
\subsection{Qualitative Analysis}
\label{sec:qualitative}
% ---------------------------------------------------------------

\chattodo{Include 3--4 bird's-eye trajectory plots from the evaluation.  The eval run already produced \texttt{videos/scene\_factory\_policy\_eval.mp4} and per-step JSONL logs that can be rendered.  Show: (1) a world where all agents succeed with clean lane-following, (2) a dense world with successful multi-agent coordination, (3) a failure case (e.g., the worst-performing world with all collisions).  Annotate start positions (green), goals (yellow), and agent trajectories (colored lines).}

Figure~\ref{fig:qualitative_trajectories} shows bird's-eye trajectory visualizations from three representative test worlds.
\chattodo{Create and include the figure.}

\paragraph{Success modes.}
In the majority of test worlds, agents follow lane centers, decelerate smoothly near goals, and avoid collisions with neighboring vehicles.
The policy learns to yield at merge points when a neighbor vehicle occupies the target lane, demonstrating emergent multi-agent coordination despite using only local observations (24 nearest neighbors within the observation radius).

\paragraph{Failure modes.}
We identify three recurring failure patterns:
\begin{enumerate}[leftmargin=*, itemsep=2pt]
  \item \textbf{Overshoot at sharp turns.}
        In worlds with tight intersections or U-turns, agents carry too much speed into the turn, overshoot the lane boundary, and trigger an off-road crash.
        This failure is absent in kinematic simulators where instantaneous heading changes are possible.
  \item \textbf{Goal approach oscillation.}
        Some agents oscillate near the goal without entering the 3.0\,m acceptance radius.
        The sign-memory brake mechanism (\S\ref{sec:action_space}) mitigates this compared to earlier training runs, but a small fraction of agents still fail to decelerate early enough.
  \item \textbf{Dense-scene congestion.}
        In worlds with all 16 agent slots occupied, vehicles occasionally deadlock when two agents attempt to merge into the same lane segment simultaneously.
        The constant-velocity TTC proxy (\S\ref{sec:sharp_edges}) provides limited lookahead for such scenarios.
\end{enumerate}

\chattodo{Verify these failure modes against the per-step JSONL logs from the eval run.  Quantify: what fraction of failures falls into each category?}
